% Meeting minutes appendix fragment.
% This file is intended to be included from midterm_report_english_sample_style.tex.

\subsection*{Minutes of the 1st Project Meeting}

\begin{tabular}{@{}p{2.4cm}p{0.72\linewidth}@{}}
  \textbf{Date} & Monday, April 20, 2026 \\
  \textbf{Time} & 2:00 pm \\
  \textbf{Place} & Online Zoom meeting \\
  \textbf{Present} & Luyan LIU; Hingchin CHEN; YSzelong LAM; Hao XU; Zhangxizi QIU; Hongrui XIAO; Yijia LI; Zhiming WANG; Rui ZHONG; Jieyi WANG; Tszwai FONG; Prof. Kai CHEN \\
  \textbf{Apology} & None \\
  \textbf{Note-taker} & Hingchin Chen \\
\end{tabular}

\paragraph{A\quad Approval of Minutes}
The meeting lasted approximately 60 minutes.

\paragraph{B\quad Discussion Items}
\begin{itemize}
  \item \textbf{Hingchin Chen:} Reported that the public FFGO workflow is based on a Wan2.2-I2V-A14B + LoRA setup and is not directly equivalent to the FastGen training interface. Proposed to first validate the native FastGen Wan2.2 TI2V/I2V path before attempting to connect a real FFGO teacher.
  \item \textbf{Prof. / leader response:} Agreed that the immediate priority should be a runnable distillation pipeline rather than final visual quality. Suggested separating the problem into two tracks: native Wan2.2 DMD2 feasibility and later FFGO teacher integration.
  \item \textbf{Action items:} Check the available WanI2V configuration, confirm the OpenVid WebDataset path, and prepare a DMD2 smoke run based on the existing FastGen Wan2.2 5B configuration.
\end{itemize}

\paragraph{C\quad Meeting Adjournment}
The meeting adjourned at 3:00 pm. The next discussion would focus on memory feasibility and whether the native CFG setting can run under the available GPU resources.

\newpage

\subsection*{Minutes of the 2nd Project Meeting}

\begin{tabular}{@{}p{2.4cm}p{0.72\linewidth}@{}}
  \textbf{Date} & Monday, April 27, 2026 \\
  \textbf{Time} & 2:00 pm \\
  \textbf{Place} & Online Zoom meeting \\
  \textbf{Present} & Luyan LIU; Hingchin CHEN; YSzelong LAM; Hao XU; Zhangxizi QIU; Hongrui XIAO; Yijia LI; Zhiming WANG; Rui ZHONG; Jieyi WANG; Tszwai FONG; Prof. Kai CHEN \\
  \textbf{Apology} & None \\
  \textbf{Note-taker} & Hingchin Chen \\
\end{tabular}

\paragraph{A\quad Approval of Minutes}
The meeting lasted approximately 60 minutes.

\paragraph{B\quad Discussion Items}
\begin{itemize}
  \item \textbf{Hingchin Chen:} Reported that FastGen native Wan2.2 TI2V 5B / DMD2 can enter the training loop but has strong memory pressure. The original CFG-on setting caused an OOM during the teacher guidance path, while a reduced-resource setting could pass the first real student update and save an initial checkpoint.
  \item \textbf{Prof. / leader response:} Recommended treating the result as a systems feasibility baseline. The feedback was to keep careful notes on GPU memory, checkpoint saving, and resume behavior, and to avoid presenting visual quality claims before there is a consistent evaluation protocol.
  \item \textbf{Action items:} Continue monitoring the 5000-iteration run, confirm checkpoint continuity, prepare fixed-prompt student inference after checkpoint saving, and keep a separate list of remaining gaps for FFGO A14B + LoRA teacher integration.
\end{itemize}

\paragraph{C\quad Meeting Adjournment}
The meeting adjourned at 3:00 pm. The next discussion would focus on whether a CFG-on low-resource variant could be made trainable.

\newpage

\subsection*{Minutes of the 3rd Project Meeting}

\begin{tabular}{@{}p{2.4cm}p{0.72\linewidth}@{}}
  \textbf{Date} & Monday, May 4, 2026 \\
  \textbf{Time} & 2:00 pm \\
  \textbf{Place} & Online Zoom meeting \\
  \textbf{Present} & Luyan LIU; Hingchin CHEN; YSzelong LAM; Hao XU; Zhangxizi QIU; Hongrui XIAO; Yijia LI; Zhiming WANG; Rui ZHONG; Jieyi WANG; Tszwai FONG; Prof. Kai CHEN \\
  \textbf{Apology} & None \\
  \textbf{Note-taker} & Hingchin Chen \\
\end{tabular}

\paragraph{A\quad Approval of Minutes}
The meeting lasted approximately 60 minutes.

\paragraph{B\quad Discussion Items}
\begin{itemize}
  \item \textbf{Hingchin Chen:} Summarized the low-resource configuration search: 81f + CFG on reached OOM, 49f became too short for the discriminator temporal kernel, and 65f + CFG on became the main trainable setting. Also reported that the 5000-iteration checkpoint sequence was complete and that resume training produced a later 0013000 checkpoint.
  \item \textbf{Prof. / leader response:} Suggested framing the midterm report around training feasibility, checkpoint continuity, and inference-flow archiving. The response was to keep visual samples as supporting assets only and leave subjective or quantitative video-quality assessment for a later stage.
  \item \textbf{Action items:} Prepare the midterm report, organize checkpoint and inference assets, list unresolved items such as the 0022000 save failure and exact quality evaluation plan, and keep future comparison with teammate configurations as an optional follow-up.
\end{itemize}

\paragraph{C\quad Meeting Adjournment}
The meeting adjourned at 3:00 pm. Follow-up work would focus on report preparation and a more systematic evaluation plan.
